% Deep Theory A↔B (Ch 24)
% Extracted from main.tex
% This file is for agent editing — will be merged back

\chapter{Deep Theory: Algorithm A $\leftrightarrow$ Algorithm B}
\label{ch:algo-deep}

Dichotomy D20 is the deepest and most challenging: two \emph{genuinely
different algorithms} that compute the same mathematical function.
Unlike D1--D19 and D21--D24, which can be resolved by syntactic
transformation or algebraic axioms, D20 requires reasoning about
the \emph{mathematical specification} that both algorithms satisfy.

This chapter develops the full CTTT theory of specification-level
equivalence, gives detailed proof sketches for every major
algorithm-equivalence class in our 100-pair benchmark, establishes
a formal boundary between structurally provable and fundamentally
extensional equivalences, and connects the theory to matroids and
the Yoneda embedding.

\section{The Problem}

Consider two programs:
\begin{itemize}
  \item $f$: Graham scan for convex hull.
  \item $g$: Andrew's monotone chain for convex hull.
\end{itemize}

These share no structural similarity: $f$ sorts by polar angle and
uses a stack with cross-product tests; $g$ sorts lexicographically
and builds upper/lower hulls.  No syntactic transformation, fold
fusion, or commutativity axiom connects them.

Yet both compute the \emph{same} mathematical function: the convex
hull of a point set.  The equivalence is a \emph{mathematical theorem}
(``the convex hull of a finite point set is unique''), not a
syntactic identity.

\section{Specifications as Propositional Types}

In CTTT, a \emph{specification} is a dependent type:

\begin{definition}[Specification]
A specification for a function $h : A \to B$ is a type family
$S : A \times B \to \Type$ such that:
\begin{enumerate}
  \item \textbf{Satisfiability}: For every $x : A$, there exists
    $y : B$ with $S(x, y)$ inhabited.
  \item \textbf{Functionality}: For every $x : A$, the type
    $\Sigma_{y:B} S(x, y)$ is \emph{propositional} (at most one
    $y$ satisfies $S(x, -)$).
\end{enumerate}
A specification satisfying both conditions is called
\emph{deterministic}.
\end{definition}

\begin{theorem}[Deterministic Specification Uniqueness]
\label{thm:det-spec}
If $S$ is a deterministic specification and both $f$ and $g$ satisfy
$S$ (i.e., $S(x, f(x))$ and $S(x, g(x))$ are inhabited for all $x$),
then $f \equiv g$.
\end{theorem}

\begin{proof}
By functionality, $\Sigma_{y:B} S(x, y)$ is propositional.  Both
$(f(x), \mathsf{pf}_f)$ and $(g(x), \mathsf{pf}_g)$ inhabit this
type.  By propositionality, they are connected by a path.  The first
projection of this path gives $\Path_B(f(x), g(x))$.  By funext
(Theorem~\ref{thm:funext}), $f \equiv g$.
\end{proof}

This is the \emph{fundamental theorem} for D20: if we can show
both algorithms satisfy the same deterministic specification,
they must compute the same function.

\section{Examples of Deterministic Specifications}

\subsection{Convex Hull}

\begin{definition}[Convex Hull Specification]
$S_{\mathrm{CH}}(P, H)$ holds iff:
\begin{enumerate}
  \item $H \subseteq P$ (the hull vertices are input points).
  \item $H$ is in convex position (no point is inside the convex
    hull of the others).
  \item Every $p \in P$ is inside or on the boundary of the convex
    hull of $H$.
  \item $H$ is listed in counterclockwise order starting from the
    leftmost-lowest point.
\end{enumerate}
\end{definition}

\begin{proposition}
$S_{\mathrm{CH}}$ is deterministic.
\end{proposition}

\begin{proof}
Uniqueness: the set of extreme points of a finite point set is unique.
The counterclockwise ordering from the leftmost-lowest point is unique.
\end{proof}

Both Graham scan and monotone chain satisfy $S_{\mathrm{CH}}$.
By Theorem~\ref{thm:det-spec}, they are equivalent.

\subsection{Shortest Paths}

\begin{definition}[Shortest Path Specification]
$S_{\mathrm{SP}}(G, s, D)$ holds iff for each vertex $v$:
\[
  D[v] = \min\{w(\pi) : \pi \text{ is a path from } s \text{ to } v
    \text{ in } G\}
\]
where $w(\pi) = \sum_{e \in \pi} w(e)$ is the path weight.
\end{definition}

\begin{proposition}
$S_{\mathrm{SP}}$ is deterministic (assuming non-negative weights
and well-defined minimum).
\end{proposition}

Both Dijkstra's algorithm and Bellman--Ford satisfy $S_{\mathrm{SP}}$
(for non-negative weights).  By Theorem~\ref{thm:det-spec}, they
produce the same distance array.

\subsection{Strongly Connected Components}

\begin{definition}[SCC Specification]
$S_{\mathrm{SCC}}(G, C)$ holds iff $C$ is the unique partition of
$V(G)$ into maximal strongly connected subsets, listed in reverse
topological order of the DAG of SCCs.
\end{definition}

Tarjan's algorithm and Kosaraju's algorithm both satisfy
$S_{\mathrm{SCC}}$, with the SCCs in the same canonical order.

\subsection{Edit Distance}

\begin{definition}[Edit Distance Specification]
$S_{\mathrm{ED}}(s_1, s_2, d)$ holds iff $d$ equals the minimum
number of single-character insertions, deletions, and substitutions
to transform $s_1$ into $s_2$.
\end{definition}

This is deterministic (the minimum exists and is unique).
Wagner--Fischer and Hirschberg both satisfy it.

\subsection{Binomial Coefficients}

\begin{definition}[Binomial Specification]
$S_{\binom{n}{k}}(n, k, v)$ holds iff $v = \frac{n!}{k!(n-k)!}$.
\end{definition}

The multiplicative formula and Pascal's triangle both compute this.
Uniqueness follows from the uniqueness of $n!$.

\section{The Specification as a Higher Inductive Type}

Each specification $S$ defines a HIT $T_S$:

\begin{definition}[Specification HIT]
Given a deterministic specification $S : A \times B \to \Type$,
the \emph{specification HIT} $T_S$ has:
\begin{itemize}
  \item \textbf{Point constructor}: for each program $f$ satisfying
    $S$, a point $\iota(f) : T_S$.
  \item \textbf{Path constructor}: for any two programs $f, g$
    satisfying $S$:
    \[
      \mathsf{det}_S : \Path_{T_S}(\iota(f), \iota(g))
    \]
  \item \textbf{Truncation}: $T_S$ is 0-truncated (propositional).
\end{itemize}
\end{definition}

The path constructor $\mathsf{det}_S$ is the computational content
of Theorem~\ref{thm:det-spec}: any two implementations of a
deterministic specification are connected.

\section{Z3 Encoding of Specifications}

A specification $S(x, y)$ is encoded as a Z3 formula over the
shared PyObj theory:

\begin{lstlisting}[caption={Specification encoding in Z3}]
# Specification: "y is the sorted permutation of x"
def sorted_spec(x, y, solver, T):
    S = T.S
    sf = T.shared_fn('sorted', 1)
    # y == sorted(x)
    solver.add(y == sf(x))
    # sorted(sorted(x)) == sorted(x) (idempotence)
    solver.add(ForAll([x], sf(sf(x)) == sf(x)))
\end{lstlisting}

\begin{lstlisting}[caption={Specification uniqueness check}]
# If both f and g satisfy sorted_spec, they are equal
x = Const('x', S)
solver.add(T.shared_fn('sorted', 1)(x) !=
           T.shared_fn('sorted', 1)(x))
# Trivially UNSAT: same shared function symbol
\end{lstlisting}

For more complex specifications (shortest paths, edit distance),
the specification is encoded as a Z3 RecFunction satisfying the
recurrence:

\begin{lstlisting}[caption={Edit distance specification}]
# d(i,j) = min(d(i-1,j)+1, d(i,j-1)+1,
#              d(i-1,j-1) + cost(i,j))
# with d(0,j)=j and d(i,0)=i
# Both DP and memoized versions satisfy this.
\end{lstlisting}

%% ═══════════════════════════════════════════════════════════
%% NEW: Detailed Proof Sketches for Algorithm Equivalence Classes
%% ═══════════════════════════════════════════════════════════

\section{Factorial Variants: The D1 Path}
\label{sec:factorial-variants}

The benchmark contains several factorial-adjacent equivalences
(EQ-21 prime factorization, EQ-48 trailing zeroes in $n!$,
EQ-40 binomial coefficients) whose proofs all pass through D1
(recursive $\leftrightarrow$ iterative).  We give the paradigmatic
proof for factorial itself, since every variant follows the same
template.

\begin{definition}[Factorial Specification]
$S_{!}(n, v)$ holds iff $n \in \Nat$ with $n \ge 0$ and
$v = \prod_{i=1}^{n} i$.  When $n = 0$, $v = 1$ (empty product).
\end{definition}

\begin{proposition}
$S_{!}$ is deterministic: the product $\prod_{i=1}^{n} i$ is a
unique integer for each $n$.
\end{proposition}

\noindent\textbf{Three implementations.}
\begin{enumerate}
  \item \emph{Recursive}:
    $\mathsf{fix}[h]\bigl(\mathsf{case}(n = 0,\; 1,\;
      n \cdot h(n-1))\bigr)$.
  \item \emph{Iterative (for-loop)}:
    $\fold[\mathsf{mul}](1, \mathsf{range}(1, n{+}1))$.
  \item \emph{functools.reduce}:
    $\mathsf{reduce}(\mathsf{operator.mul},\;
      \mathsf{range}(1, n{+}1),\; 1)$.
\end{enumerate}

\begin{theorem}[Factorial Variant Equivalence]
\label{thm:factorial-equiv}
All three implementations satisfy $S_{!}$ and are therefore
equal by Theorem~\ref{thm:det-spec}.  The structural path is:
\[
  \mathsf{fix}[h](\ldots) \xrightarrow{\text{D1\_fix\_to\_fold}}
  \fold[\mathsf{mul}](1, \mathsf{range})
  \xleftarrow{\text{D5\_fold\_universal}}
  \mathsf{reduce}(\mathsf{mul}, \mathsf{range}, 1)
\]
\end{theorem}

\begin{proof}[Proof sketch]
\textbf{Step 1} (Recursive $\to$ Fold).  The recursive body has the
shape $h(n) = n \cdot h(n-1)$ with base case $h(0) = 1$.  This is a
\emph{primitive recursion} on $\Nat$ with accumulation operator
$\mathsf{mul}$ and initial value $1$.  By the D1 axiom
(\texttt{D1\_fix\_to\_fold} in \texttt{path\_search.py}), we extract:
\[
  \mathsf{fix}[h](\ldots) \;\equiv\;
  \fold[\mathsf{mul}]\bigl(1, \mathsf{range}(1, n{+}1)\bigr)
\]
The OTerm normalizer (phase~3: fold recognition) performs this
automatically by detecting that the fix-point body is a simple
accumulation over a decreasing integer argument.

\textbf{Step 2} (reduce $\to$ Fold).  The \texttt{functools.reduce}
call with $(\mathsf{operator.mul}, \mathsf{range}(1,n{+}1), 1)$
is compiled by the denotational semantics (\texttt{denotation.py},
phase~4: HOF fusion) into $\fold[\mathsf{mul}](1, \mathsf{range}(1,n{+}1))$
--- the same OTerm.  The D5 axiom (\texttt{D5\_fold\_universal})
confirms fold universality: different surface representations of the
same fold (explicit loop, \texttt{reduce}, comprehension) all compile
to $\fold[\mathsf{op}](\mathsf{init}, \mathsf{coll})$.

\textbf{Step 3} (Canonical form).  All three implementations share
the canonical OTerm:
\[
  \fold[\mathsf{mul}](1, \mathsf{range}(1, \mathsf{add}(n, 1)))
\]
By \texttt{refl}, they are identical.  Alternatively, spec
identification (\texttt{\_identify\_spec} in \texttt{path\_search.py})
recognizes $\fold[\mathsf{mul}](1, \mathsf{range}(\ldots))$ as the
``factorial'' specification, lifting all three to
$\mathsf{@factorial}(n)$ via D20.
\end{proof}

\begin{remark}[Application to EQ-40: Binomial Coefficients]
The multiplicative formula
$\binom{n}{k} = \prod_{i=0}^{k-1} \frac{n-i}{i+1}$
compiles to
$\fold[\mathsf{div\_mul}](1, \mathsf{range}(k))$.
Pascal's triangle compiles to a nested fold over rows.
Both satisfy $S_{\binom{n}{k}}$.  Since $\binom{n}{k}$ is unique,
Theorem~\ref{thm:det-spec} gives the path.  This is an instance of
D1 (fold recognition) plus D20 (spec uniqueness), because the two
folds have different recurrence structures that happen to compute
the same arithmetic function.
\end{remark}

\section{Fibonacci Variants: D1 and D16 Paths}
\label{sec:fibonacci-variants}

EQ-37 (matrix exponentiation vs.\ fast doubling for Fibonacci mod $p$)
is one of the deepest structural equivalences in the benchmark.  The
two algorithms share no syntactic structure, yet both compute
$F(n) \bmod p$.

\begin{definition}[Fibonacci Specification]
$S_{\mathrm{fib}}(n, v)$ holds iff $v = F(n)$ where $F$ is the unique
function satisfying $F(0) = 0$, $F(1) = 1$,
$F(n) = F(n{-}1) + F(n{-}2)$ for $n \ge 2$.
\end{definition}

\noindent\textbf{Five implementations.}
\begin{enumerate}
  \item \emph{Naive recursive}:
    $\mathsf{fix}[f]\bigl(\mathsf{case}(n \le 1, n,
      f(n{-}1) + f(n{-}2))\bigr)$.
  \item \emph{Memoized recursive} (top-down DP): same recurrence,
    with a memo table wrapping the fix-point.
  \item \emph{Iterative} (bottom-up DP):
    $\fold[\mathsf{fib\_step}]((0,1), \mathsf{range}(n))$ where
    $\mathsf{fib\_step}((\mathit{a},\mathit{b}), \_) =
      (\mathit{b}, \mathit{a}+\mathit{b})$.
  \item \emph{Matrix exponentiation} (EQ-37a):
    $\bigl(\begin{smallmatrix}1&1\\1&0\end{smallmatrix}\bigr)^n$
    computed by repeated squaring; result is entry $[0][1]$.
  \item \emph{Fast doubling} (EQ-37b): uses the identities
    $F(2k) = F(k)[2F(k{+}1) - F(k)]$ and
    $F(2k{+}1) = F(k)^2 + F(k{+}1)^2$.
\end{enumerate}

\begin{theorem}[Fibonacci Variant Equivalence]
\label{thm:fib-equiv}
All five implementations satisfy $S_{\mathrm{fib}}$.  The proof paths are:
\begin{itemize}
  \item \emph{Naive $\leftrightarrow$ Memoized}: D16 (memo $\leftrightarrow$ DP).
    Both share the same recurrence
    $f(n) = f(n{-}1) + f(n{-}2)$.  Memoization only changes
    \emph{evaluation order}, not the mathematical function.  The
    D16 axiom (\texttt{D16\_recurrence\_normalize}) normalizes
    both to the same canonical fix-point.
  \item \emph{Memoized $\leftrightarrow$ Iterative}: D1
    (fix $\to$ fold).  The fix-point with memo has a linear
    dependency structure (each $f(k)$ depends only on $f(k{-}1)$
    and $f(k{-}2)$), so it compiles to a fold with state
    $(a, b)$.
  \item \emph{Iterative $\leftrightarrow$ Matrix}: D20
    (spec uniqueness).  The fold computes $F(n)$ by the
    recurrence; matrix exponentiation computes $F(n)$ by the
    identity $\bigl(\begin{smallmatrix}1&1\\1&0\end{smallmatrix}\bigr)^n
    = \bigl(\begin{smallmatrix}F(n{+}1)&F(n)\\
    F(n)&F(n{-}1)\end{smallmatrix}\bigr)$.
    Both satisfy $S_{\mathrm{fib}}$; uniqueness of $F(n)$ gives
    the path.
  \item \emph{Matrix $\leftrightarrow$ Fast doubling}: D20
    (spec uniqueness).  The fast-doubling identities are
    consequences of the matrix identity above.  Both satisfy
    $S_{\mathrm{fib}}$.
\end{itemize}
\end{theorem}

\begin{proof}[Proof sketch for Matrix $\leftrightarrow$ Fast Doubling]
The key is that both algorithms implement efficient computation of
the same recurrence by exploiting the multiplicative structure of
$\Nat$.  Matrix exponentiation uses
$M^{2k} = (M^k)^2$ and $M^{2k+1} = M \cdot (M^k)^2$.
Fast doubling derives closed-form identities from this matrix
equation.  Define $M \coloneqq
\bigl(\begin{smallmatrix}1&1\\1&0\end{smallmatrix}\bigr)$.
Then:
\[
  M^k = \begin{pmatrix} F(k{+}1) & F(k) \\ F(k) & F(k{-}1)
  \end{pmatrix}
\]
Squaring: $M^{2k} = (M^k)^2$ gives, by entry-wise expansion:
\begin{align*}
  F(2k) &= F(k)\bigl[2F(k{+}1) - F(k)\bigr] \\
  F(2k{+}1) &= F(k)^2 + F(k{+}1)^2
\end{align*}
These are exactly the fast-doubling identities.  Both compute
$F(n)$ in $O(\log n)$ multiplications.  Since $S_{\mathrm{fib}}$
is deterministic, the path exists by Theorem~\ref{thm:det-spec}.

In OTerm form, both compile to $\mathsf{@fibonacci\_like}(n)$ via
\texttt{\_identify\_spec}, which recognizes any fix-point whose
body contains additions and subtractions in the Fibonacci pattern.
\end{proof}

\section{Sorting-Then-Processing: The D19 Path}
\label{sec:sort-process}

Many benchmark pairs share the pattern: \emph{sort the data, then
process it}.  Different sort strategies (key function, reverse flag,
built-in \texttt{sorted} vs.\ manual merge sort) with the same
post-processing produce identical outputs.

\begin{definition}[Sort-Process Specification]
Let $\sigma : \mathcal{P}(A) \to \mathsf{List}(A)$ be a sorting
function and $\phi : \mathsf{List}(A) \to B$ a processing function.
The specification is $S(x, y)$ iff $y = \phi(\sigma(x))$.

When $\sigma$ is determined by a total order on $A$ (possibly
with a key function), $S$ is deterministic: the sorted order is
unique and $\phi$ is a function.
\end{definition}

\noindent\textbf{Benchmark instances.}
\begin{itemize}
  \item \textbf{EQ-01} (flatten + sort + unique): both
    implementations recursively flatten, then apply
    $\mathsf{sorted}(\mathsf{set}(\ldots))$.  The structural
    path is D4 (comprehension fusion) $+$ D19 (sort canonicalization):
    different flattening strategies (recursive generator vs.\
    iterative stack) both produce the same multiset, and
    $\mathsf{sorted} \circ \mathsf{set}$ is a canonical
    representative of the equivalence class.
  \item \textbf{EQ-02} (word frequency histogram): both produce
    $\mathsf{sorted}(\mathsf{Counter}(\mathsf{words}).
    \mathsf{items}())$.  The D13 axiom (Counter $\leftrightarrow$
    manual fold) and D19 (sort absorption) give the path.
  \item \textbf{EQ-25} (merge overlapping intervals): both sort
    intervals by start time, then sweep.  The path is D19:
    $\fold[\mathsf{merge}](\mathsf{init},
    \mathsf{sorted}(\mathsf{intervals}))$
    is the canonical form for both.
  \item \textbf{EQ-38} (group anagrams): one uses
    $\mathsf{sorted}(\mathsf{word})$ as grouping key; the other
    uses a prime-product hash.  Both produce the same partition
    of strings into anagram classes.  The path requires D20
    (spec uniqueness): the anagram-partition specification is
    deterministic when the output is canonicalized by sorting
    both within groups and across groups.
\end{itemize}

\begin{proposition}[Sort Absorption]
\label{prop:sort-absorb}
In the OTerm algebra, $\mathsf{sorted}(\mathsf{sorted}(x)) \equiv
\mathsf{sorted}(x)$ (idempotence).  Furthermore, for any
commutative and associative fold operation $\oplus$:
\[
  \fold[\oplus](\mathsf{init}, \mathsf{sorted}(x)) \equiv
  \fold[\oplus](\mathsf{init}, x)
\]
because $\oplus$ is order-invariant.  The D19 axiom
(\texttt{D19\_sort\_irrelevant}) implements this.
\end{proposition}

\section{GCD Variants: The D1 Path with Recurrence Matching}
\label{sec:gcd-variants}

\begin{definition}[GCD Specification]
$S_{\gcd}(a, b, d)$ holds iff $d$ is the greatest common divisor:
$d \mid a$, $d \mid b$, and for all $c$ with $c \mid a$ and
$c \mid b$, $c \mid d$.
\end{definition}

\begin{proposition}
$S_{\gcd}$ is deterministic: $\gcd(a, b)$ is unique for $a, b > 0$.
\end{proposition}

\noindent\textbf{Two implementations.}
\begin{enumerate}
  \item \emph{Euclidean} (division-based):
    $\mathsf{fix}[g]\bigl(\mathsf{case}(b = 0,\; a,\;
      g(b, a \bmod b))\bigr)$.
  \item \emph{Subtraction-based}:
    $\mathsf{fix}[g]\bigl(\mathsf{case}(a = b,\; a,\;
      \mathsf{case}(a > b,\; g(a{-}b, b),\; g(a, b{-}a)))\bigr)$.
\end{enumerate}

\begin{theorem}[GCD Variant Equivalence]
Both implementations satisfy $S_{\gcd}$.  The proof uses D1
(fix-point equivalence) with \emph{recurrence matching}:
both compute the same function despite different reduction strategies.
\end{theorem}

\begin{proof}[Proof sketch]
The Euclidean algorithm reduces $(a, b) \mapsto (b, a \bmod b)$;
the subtraction algorithm reduces
$(a, b) \mapsto (a - b, b)$ or $(a, b - a)$.
Both terminate (the sum $a + b$ strictly decreases in the
Euclidean case; the larger argument strictly decreases in the
subtraction case).  At termination, both return the last nonzero
value, which equals $\gcd(a, b)$ by the fundamental property
$\gcd(a, b) = \gcd(b, a \bmod b) = \gcd(a - b, b)$.

In OTerm form, both compile to $\mathsf{fix}$ nodes.  The D1
path does not directly connect them (the fix-point bodies have
different structure).  Instead, D20 applies: both satisfy the
deterministic specification $S_{\gcd}$, so by
Theorem~\ref{thm:det-spec}, they are equal.

Note that the subtraction-based algorithm is strictly less
efficient ($O(\max(a,b))$ vs.\ $O(\log \min(a,b))$), but
efficiency is invisible in the extensional type theory: the
path connects their \emph{values}, not their computation traces.
\end{proof}

%% ═══════════════════════════════════════════════════════════
%% The Yoneda Perspective (expanded)
%% ═══════════════════════════════════════════════════════════

\section{Algorithm Equivalence via Yoneda}
\label{sec:yoneda-algo}

The Yoneda lemma provides a deep perspective on specification-level
equivalence that goes beyond the ad hoc ``identify a specification''
approach.

\begin{theorem}[Yoneda Lemma]
For a presheaf $F : \mathcal{C}^{\mathrm{op}} \to \Set$ and an
object $A \in \mathcal{C}$:
\[
  \mathrm{Nat}(\mathrm{Hom}(A, -), F) \cong F(A)
\]
\end{theorem}

\noindent\textbf{Interpretation for programs.}  Let $\mathcal{C}$
be the category whose objects are Python types and whose morphisms
are Python functions.  A program $f : A \to B$ defines a
representable presheaf $\mathrm{Hom}(B, -)$.  The Yoneda lemma
says: $f$ is \emph{uniquely determined} by how it interacts with
\emph{all} observations $t : B \to C$:
\[
  f \equiv g \iff \forall C, \forall t : B \to C.\;
    t \circ f \equiv t \circ g
\]

This is the \emph{observational equivalence} formulation: programs
are characterized by their specifications, not their implementations.
Two programs are the same iff no downstream consumer can tell them
apart.

\begin{corollary}[Yoneda Embedding for Programs]
\label{cor:yoneda-embed}
The Yoneda embedding $\mathsf{y} : \mathcal{C} \to [\mathcal{C}^{\mathrm{op}}, \Set]$
is full and faithful.  Applied to programs: two implementations
$f, g : A \to B$ are equal in $\mathcal{C}$ iff their Yoneda
images are naturally isomorphic.  That is, a program \emph{is}
its specification.
\end{corollary}

\noindent\textbf{Z3 implementation.}
In Z3, observations are approximated by
the shared function symbols $t \in \{$\texttt{sorted}, \texttt{len},
\texttt{sum}, \texttt{hash}, \texttt{str}, $\ldots\}$.  Two programs
$f, g$ are observationally equivalent iff for every shared symbol
$t$:
\[
  \forall p.\; t(\den{f}(p)) = t(\den{g}(p))
\]
This is checked as a single Z3 query: $\exists p.\, \bigvee_t
t(\den{f}(p)) \neq t(\den{g}(p))$.  If UNSAT, the programs are
Yoneda-equivalent up to the finite observation set.

\begin{remark}[Finite Approximation]
The shared function symbols form a \emph{finite} approximation of
the full Yoneda embedding.  Soundness: if the Z3 query returns
UNSAT, the programs agree on all modeled observations.
Completeness gap: there may exist an observation not modeled by
any shared symbol.  This gap is narrowed by adding more shared
symbols (the theory library grows over time) and is bridged entirely
for specifications expressible in the Z3 theory.
\end{remark}

%% ═══════════════════════════════════════════════════════════
%% Matroid / Greedy Theorem (NEW)
%% ═══════════════════════════════════════════════════════════

\section{The Matroid/Greedy Theorem and D20}
\label{sec:matroid-greedy}

A fundamental question in algorithm theory:
\emph{when does a greedy algorithm compute an optimal solution?}
The classical answer involves matroids.  We now connect this to
D20 and the CTTT framework.

\begin{definition}[Matroid]
A \emph{matroid} $M = (E, \mathcal{I})$ consists of a finite ground
set $E$ and a family $\mathcal{I} \subseteq 2^E$ of \emph{independent
sets} satisfying:
\begin{enumerate}
  \item $\emptyset \in \mathcal{I}$ (the empty set is independent).
  \item If $A \in \mathcal{I}$ and $B \subseteq A$, then
    $B \in \mathcal{I}$ (hereditary property).
  \item If $A, B \in \mathcal{I}$ and $|A| < |B|$, then there
    exists $e \in B \setminus A$ with $A \cup \{e\} \in \mathcal{I}$
    (exchange property).
\end{enumerate}
\end{definition}

\begin{theorem}[Rado--Edmonds Greedy Theorem]
\label{thm:matroid-greedy}
A greedy algorithm (iteratively adding the cheapest feasible element)
solves the weighted optimization problem
$\max_{I \in \mathcal{I}} \sum_{e \in I} w(e)$
if and only if $\mathcal{I}$ forms a matroid.
\end{theorem}

\noindent\textbf{Connection to D20.}
The matroid theorem provides a \emph{metatheorem} for equivalence:
if the feasibility system forms a matroid, then the greedy algorithm
and any DP/brute-force algorithm computing the optimal solution
satisfy the same deterministic specification $S_{\mathrm{opt}}(x, v)
\coloneqq v = \max_{I \in \mathcal{I}} \sum_{e \in I} w(e)$.

\begin{corollary}[Greedy $\leftrightarrow$ DP on Matroids]
\label{cor:greedy-dp}
If the optimization problem has matroid structure, then the greedy
algorithm and the DP algorithm are D20-equivalent: both satisfy the
deterministic specification $S_{\mathrm{opt}}$, so
Theorem~\ref{thm:det-spec} provides the path.
\end{corollary}

\noindent\textbf{Benchmark instance: EQ-47} (stock profit with $k$
transactions).  The DP algorithm (\texttt{eq47a}) fills a 2D table;
the alternative (\texttt{eq47b}) reduces to a greedy sum when
$k \ge n/2$ and falls back to DP otherwise.  The specification is:
\[
  S_{\mathrm{profit}}(\mathsf{prices}, k, v) \iff
  v = \max\left\{\sum_{j=1}^{m}
    (\mathsf{prices}[s_j] - \mathsf{prices}[b_j])
    : m \le k,\; b_1 < s_1 < b_2 < s_2 < \cdots\right\}
\]
The maximum is unique (it is a real number), so $S_{\mathrm{profit}}$
is deterministic.  When $k \ge n/2$, every profitable adjacent-day
trade is taken --- this is the greedy solution on the ``pick
non-overlapping profit intervals'' matroid.  The D20 path connects
the greedy branch to the DP solution.

The D18 axiom (\texttt{D18\_greedy} in \texttt{path\_search.py}) is
the structural counterpart: it fires when the path search detects
that a greedy algorithm and a DP algorithm operate on the same
optimization problem with matroid structure.

%% ═══════════════════════════════════════════════════════════
%% Constructive Witness for D20 (retained + expanded)
%% ═══════════════════════════════════════════════════════════

\section{Constructive Witness for D20}

For each D20 pair, the equivalence proof has the structure:

\begin{enumerate}
  \item \textbf{Identify the specification} $S$ that both programs
    satisfy.
  \item \textbf{Prove $S$ is deterministic}: show that
    $\Sigma_y S(x, y)$ is propositional for all $x$.
  \item \textbf{Construct the path}: apply Theorem~\ref{thm:det-spec}
    to obtain $f \equiv g$.
\end{enumerate}

Step 1 is done by the Z3 semantic algebra: the shared function
symbols and axioms encode the specification implicitly.  When both
programs use the same shared functions with the same axioms, they
inhabit the same specification.  The \texttt{\_identify\_spec}
function in \texttt{path\_search.py} automates this for common
specifications (factorial, Fibonacci, sorted, binomial).

Step 2 is done by the Z3 axioms: the idempotence, commutativity,
and associativity axioms ensure that specifications involving sorting,
folding, etc., are deterministic.

Step 3 is done by Theorem~\ref{thm:det-spec} (which is a metatheorem
of CTTT, not something Z3 needs to prove).

\section{Limits of D20}

D20 cannot resolve all algorithm-level equivalences.  It fails when:
\begin{itemize}
  \item The specification is not deterministic (e.g., ``find any
    spanning tree'' --- there may be multiple valid outputs).
  \item The specification cannot be expressed in the Z3 theory
    (e.g., requires reasoning about infinite structures).
  \item The programs satisfy different specifications that happen
    to agree on all finite inputs but differ on edge cases.
\end{itemize}

We now develop the full theory for non-deterministic specifications.

\section{Quotient Types for Non-Deterministic Specifications}

A \emph{non-deterministic} specification allows multiple valid outputs
for a single input.  Examples:
\begin{itemize}
  \item ``Find \emph{any} spanning tree of the graph.''
  \item ``Return \emph{a} topological ordering'' (there may be many).
  \item ``Find \emph{a} shortest path'' (there may be ties).
  \item ``Return \emph{any} element satisfying the predicate.''
\end{itemize}

For such specifications, Theorem~\ref{thm:det-spec} does not apply:
two programs can both satisfy the specification yet produce different
(equally valid) outputs.  We need a richer type-theoretic structure.

\begin{definition}[Quotient Type]
\label{def:quotient}
Given a type $B$ and an equivalence relation $\sim$ on $B$, the
\emph{quotient type} $B / {\sim}$ is the HIT with:
\begin{itemize}
  \item Point constructor: $q : B \to B / {\sim}$
    (the quotient map).
  \item Path constructor: for all $b_1, b_2 : B$ with $b_1 \sim b_2$:
    \[
      \mathsf{quot} : \Path_{B/{\sim}}(q(b_1), q(b_2))
    \]
  \item Truncation: $B / {\sim}$ is a set (0-truncated).
\end{itemize}
\end{definition}

\begin{definition}[Non-Deterministic Specification via Quotient]
\label{def:nondet-spec}
A non-deterministic specification $S : A \to \mathcal{P}(B)$ (mapping
each input to a \emph{set} of valid outputs) defines an equivalence
relation on outputs:
\[
  b_1 \sim_{S,x} b_2 \iff b_1 \in S(x) \land b_2 \in S(x)
\]
(two outputs are equivalent if they are both valid for input $x$).

Two programs $f, g$ are $S$-equivalent iff for all $x$:
\[
  q(f(x)) = q(g(x)) \quad \text{in } B / {\sim_{S,x}}
\]
That is, their outputs are in the same equivalence class --- both
are valid responses to the specification.
\end{definition}

\begin{theorem}[Non-Deterministic Specification Equivalence]
\label{thm:nondet-equiv}
Two programs $f, g : A \to B$ are $S$-equivalent iff for all $x : A$:
\[
  f(x) \in S(x) \;\land\; g(x) \in S(x)
\]
The cubical path witnessing this equivalence is:
\[
  p_x = \mathsf{quot}(f(x), g(x)) : \Path_{B/{\sim_{S,x}}}(q(f(x)), q(g(x)))
\]
which exists because both $f(x)$ and $g(x)$ are in $S(x)$, so
$f(x) \sim_{S,x} g(x)$.
\end{theorem}

\begin{proof}
By Definition~\ref{def:quotient}, if $f(x) \sim_{S,x} g(x)$ (both
are valid), then $\mathsf{quot}$ provides a path in $B / {\sim_{S,x}}$.
By function extensionality, the family $\lambda x.\, p_x$ gives:
\[
  p : \Path_{A \to B/{\sim_S}}(q \circ f, \; q \circ g)
\]
\end{proof}

\subsection{Z3 Encoding of Quotient Types}

The quotient is encoded in Z3 by replacing the inequality query.
Instead of asking ``$\exists x.\, f(x) \neq g(x)$'', we ask:

\[
  \exists x.\; f(x) \notin S(x) \;\lor\; g(x) \notin S(x)
\]

If UNSAT: both programs satisfy $S$ on all inputs, so they are
$S$-equivalent (even if $f(x) \neq g(x)$ for some $x$).

If SAT: at least one program violates $S$ --- the counterexample
identifies which program fails and on what input.

\begin{lstlisting}[caption={Non-deterministic spec check in Z3}]
def check_nondet_equiv(f_term, g_term, spec_pred, params, solver):
    """Check S-equivalence: both f and g satisfy spec S.

    spec_pred(x, y) is a Z3 BoolRef that is True iff y in S(x).
    """
    x = params[0]  # symbolic input

    # Check: does f violate S?
    solver.push()
    solver.add(Not(spec_pred(x, f_term)))
    f_violates = solver.check()
    solver.pop()

    # Check: does g violate S?
    solver.push()
    solver.add(Not(spec_pred(x, g_term)))
    g_violates = solver.check()
    solver.pop()

    if f_violates == unsat and g_violates == unsat:
        return EQUIVALENT  # both satisfy S everywhere
    elif f_violates == sat:
        return BUG_IN_F, solver.model()
    elif g_violates == sat:
        return BUG_IN_G, solver.model()
    else:
        return UNKNOWN
\end{lstlisting}

\subsection{Example: Spanning Tree}

\begin{lstlisting}[caption={Two spanning tree algorithms}]
def kruskal(graph):
    """Kruskal's MST (sorted by weight, union-find)."""
    edges = sorted(graph.edges, key=lambda e: e.weight)
    # ... union-find MST construction ...
    return tree

def prim(graph):
    """Prim's MST (priority queue from start vertex)."""
    # ... heap-based MST construction ...
    return tree
\end{lstlisting}

The specification $S_{\mathrm{MST}}(G, T)$ is:
\begin{enumerate}
  \item $T$ is a spanning tree of $G$ (connected, acyclic, spans all vertices).
  \item $T$ has minimum total weight among all spanning trees.
\end{enumerate}

When edge weights are distinct, the MST is unique and $S$ is
deterministic.  When weights have ties, multiple valid MSTs exist.

In the non-deterministic case:
\begin{itemize}
  \item $\mathsf{kruskal}(G)$ and $\mathsf{prim}(G)$ may produce
    \emph{different} MSTs (due to tie-breaking).
  \item But both satisfy $S_{\mathrm{MST}}$.
  \item In the quotient $B / {\sim_{S_{\mathrm{MST}}}}$, both outputs
    are in the same equivalence class.
  \item The Z3 check verifies that both satisfy $S$ (UNSAT on both
    negations), concluding $S$-equivalence.
\end{itemize}

\subsection{Example: Topological Sort}

The specification ``return a topological ordering of the DAG'' is
non-deterministic: a DAG with $n$ vertices may have exponentially
many valid orderings.

Two topological sort algorithms (Kahn's BFS vs.\ DFS-based) both
satisfy the specification:
\[
  S_{\mathrm{topo}}(G, \sigma) \iff \forall (u,v) \in E(G).\;
    \sigma^{-1}(u) < \sigma^{-1}(v)
\]
(every edge goes from earlier to later in the ordering).

They may produce different orderings, but both are valid.  In the
quotient, they are equal.

If the specification is strengthened to ``return the \emph{lexicographically
smallest} topological ordering'', it becomes deterministic, and
Theorem~\ref{thm:det-spec} applies directly.

\section{Propositional Extensionality}

An alternative to quotient types is \emph{propositional extensionality}:

\begin{axiom}[Propositional Extensionality]
\label{ax:propext}
For propositions $P, Q : \mathsf{Prop}$:
\[
  (P \leftrightarrow Q) \to (P = Q)
\]
Two propositions with the same truth value are \emph{equal}.
\end{axiom}

In cubical type theory, propositional extensionality is a consequence
of univalence (Axiom~\ref{ax:univalence}).

\textbf{Application to specifications}: if the specification $S(x, y)$
is a proposition (either inhabited or empty, with no computational
content), then propositional extensionality says:
\[
  S(x, f(x)) \leftrightarrow S(x, g(x)) \implies
  S(x, f(x)) = S(x, g(x))
\]

So if both $f(x)$ and $g(x)$ satisfy $S$ (both sides are $\top$),
they are \emph{equal as propositions}.  This doesn't make $f(x)$
and $g(x)$ equal as \emph{values}, but it makes their \emph{validity
status} equal.

Combined with the quotient construction: if validity is all we care
about (not the specific output), propositional extensionality
collapses the quotient to a single point.

\section{Set-Level Truncation}

The \emph{set-level truncation} $\|B\|_0$ of a type $B$ is the
quotient that identifies all elements:

\begin{definition}[Set Truncation]
$\|B\|_0$ is the HIT with:
\begin{itemize}
  \item $|b| : \|B\|_0$ for each $b : B$.
  \item $\mathsf{trunc} : \prod_{x, y : \|B\|_0} \Path(x, y)$
    (all elements are equal).
\end{itemize}
\end{definition}

Applied to specifications: $\|S(x)\|_0$ is either empty (no valid
output) or a single point (at least one valid output exists, and
all valid outputs are identified).

Two programs $f, g$ satisfying $S$ produce elements $|f(x)|, |g(x)|
\in \|S(x)\|_0$.  By truncation, $|f(x)| = |g(x)|$.

\begin{theorem}[Equivalence Modulo Non-Deterministic Spec]
If $S$ is a (possibly non-deterministic) specification and both $f$
and $g$ satisfy $S$ for all inputs, then:
\[
  q \circ f \equiv q \circ g : A \to \|B\|_0 / {\sim_S}
\]
where $q$ is the quotient-then-truncate map.
\end{theorem}

\begin{proof}
For each $x$, both $f(x)$ and $g(x)$ are in $S(x)$.  The quotient
$B / {\sim_{S,x}}$ identifies them (they are in the same equivalence
class).  The truncation $\|B / {\sim_{S,x}}\|_0$ ensures this is a
mere proposition.  By propositional extensionality,
$|q(f(x))| = |q(g(x))|$.  By funext, $q \circ f = q \circ g$.
\end{proof}

\section{Z3 Encoding of Truncated Specifications}

In practice, the Z3 encoding of truncated specifications is simpler
than the type theory suggests:

\begin{lstlisting}[caption={Truncated spec check}]
def check_truncated_equiv(f_term, g_term, spec_pred,
                          params, solver):
    """Check: both satisfy spec S.

    We do NOT check f(x) == g(x).
    We check f(x) in S(x) AND g(x) in S(x).
    If both hold, they are equivalent modulo S.
    """
    x = params[0]

    # Combined query: exists x where f or g violates S
    solver.push()
    solver.add(Or(
        Not(spec_pred(x, f_term)),
        Not(spec_pred(x, g_term))
    ))
    result = solver.check()
    solver.pop()

    if result == unsat:
        return S_EQUIVALENT  # both satisfy S everywhere
    else:
        m = solver.model()
        # Which one fails?
        # ... extract counterexample ...
        return S_INEQUIVALENT, m
\end{lstlisting}

\subsection{Fiber-Wise Truncated Spec Checking}

The cohomological decomposition applies to truncated specs:

\begin{enumerate}
  \item For each type fiber $U_\tau$: check that both $f$ and $g$
    satisfy $S$ when restricted to fiber $\tau$.
  \item Build the $C^0$ cochain: entry is 1 iff both satisfy $S$
    on that fiber.
  \item Compute $\Hone$: measures whether the local $S$-satisfactions
    are globally coherent.
  \item $\Hone = 0 \implies$ globally $S$-equivalent.
\end{enumerate}

The only change from deterministic spec checking is the local query:
instead of ``$f(x) = g(x)$'' we check ``$f(x) \in S(x) \land
g(x) \in S(x)$''.  The cohomological gluing machinery is unchanged.

%% ═══════════════════════════════════════════════════════════
%% The Hard 8 Analysis (NEW)
%% ═══════════════════════════════════════════════════════════

\section{The ``Hard 8'' Analysis}
\label{sec:hard-8}

Among the 50~equivalent pairs in the benchmark, most are resolved
by structural path constructors (D1--D19, D21--D24).  Eight pairs,
however, resist \emph{all} structural paths and require genuine
specification-level reasoning (D20).  We call these the ``hard~8.''

\begin{center}
\small
\begin{tabular}{r l l p{4.5cm} l}
  \textbf{ID} & \textbf{Alg A} & \textbf{Alg B} &
  \textbf{Specification} & \textbf{Why hard} \\
  \hline
  07 & Graham scan & Monotone chain &
    Convex hull &
    Different invariants \\
  08 & Wagner--Fischer & Hirschberg &
    Edit distance &
    Space-time tradeoff \\
  16 & Tarjan SCC & Kosaraju SCC &
    Strongly connected components &
    Different traversals \\
  29 & deepcopy + sort & BFS copy + sort &
    Deep copy canonical form &
    Different copy strategies \\
  33 & Factoradic & itertools.perm &
    $n$-th permutation &
    Combinatorial identity \\
  37 & Matrix exp & Fast doubling &
    $F(n) \bmod p$ &
    Algebraic identity \\
  40 & Multiplicative & Pascal's triangle &
    $\binom{n}{k}$ &
    Different recurrences \\
  47 & DP & Greedy + DP &
    Max stock profit &
    Matroid reduction \\
\end{tabular}
\end{center}

\subsection{Why Structural Paths Fail}

For each hard pair, we explain why the 23 non-D20 path constructors
cannot connect them:

\begin{enumerate}
  \item \textbf{EQ-07 (Graham $\leftrightarrow$ Monotone Chain).}
    Graham sorts by \emph{polar angle} from a pivot and maintains a
    CCW stack.  Monotone chain sorts \emph{lexicographically} and
    builds two half-hulls.  The sort keys differ, the loop invariants
    differ, and the intermediate data structures differ.  No
    combination of D1 (fold/unfold), D4 (fusion), D5 (fold
    universality), or D19 (sort-process) connects them, because the
    \emph{sort itself} is part of the algorithm's correctness argument,
    not just a preprocessing step.

  \item \textbf{EQ-08 (Wagner--Fischer $\leftrightarrow$ Hirschberg).}
    Both compute edit distance, but Hirschberg uses a
    divide-and-conquer strategy with $O(n)$ space.  The D17
    (divide-and-conquer $\leftrightarrow$ iterative) axiom applies
    only when the combining operation is associative; here, the
    ``combine'' step is a nontrivial alignment merge.

  \item \textbf{EQ-16 (Tarjan $\leftrightarrow$ Kosaraju).}
    Tarjan uses a single DFS with a low-link stack.
    Kosaraju uses two DFS passes (one on the original graph, one on
    the transpose).  The traversal structures are fundamentally
    different.  D15 (BFS $\leftrightarrow$ DFS) does not apply
    because the algorithms are not computing a fold over a traversal
    --- they use the traversal structure itself to identify SCCs.

  \item \textbf{EQ-37 (Matrix Exp $\leftrightarrow$ Fast Doubling).}
    Addressed in \S\ref{sec:fibonacci-variants}.  The matrix and
    scalar recurrences have different OTerm structures despite
    computing the same sequence.

  \item \textbf{EQ-47 (DP $\leftrightarrow$ Greedy+DP).}
    Addressed in \S\ref{sec:matroid-greedy}.  The greedy branch is
    not a structural simplification of the DP --- it uses a
    fundamentally different algorithm justified by the matroid
    property.
\end{enumerate}

\subsection{The Common Pattern}

All hard-8 pairs share a common feature: the two algorithms use
\emph{different mathematical insights} to solve the same problem.
The structural paths (D1--D19, D21--D24) transform one algorithm
into another by rewriting its implementation; the hard-8 pairs
cannot be connected this way because the implementations encode
genuinely different proofs of the same theorem.

In CTTT terms: the hard-8 pairs require paths in the
\emph{specification HIT} $T_S$, not paths in the \emph{program HIT}
$\mathsf{PyComp}$.  The path constructor $\mathsf{det}_S$
(Definition of the Specification HIT) is the only axiom that
applies, and it requires showing that both algorithms satisfy the
same deterministic specification --- a mathematical argument, not a
syntactic rewrite.

%% ═══════════════════════════════════════════════════════════
%% Formal Boundary Theorem (NEW)
%% ═══════════════════════════════════════════════════════════

\section{The Structural/Extensional Boundary}
\label{sec:boundary}

We now formalize the distinction between equivalences provable by
structural rewriting and those requiring extensional (specification-level)
reasoning.

\begin{definition}[Structural Equivalence]
\label{def:structural-equiv}
Two OTerms $t_1, t_2$ are \emph{structurally equivalent}, written
$t_1 \sim_{\mathsf{str}} t_2$, if there exists a path in the HIT
$\mathsf{PyComp}$ using only the path constructors from D1--D19
and D21--D24 (the non-D20 axioms).

Formally, $\sim_{\mathsf{str}}$ is the equivalence relation generated
by the rewrite rules of \texttt{path\_search.py}: the axiom functions
\texttt{\_axiom\_d1\_fold\_unfold} through
\texttt{\_axiom\_d24\_eta}, plus the algebraic axioms
(commutativity, associativity, identity, involution), quotient axioms,
fold axioms, and case axioms, applied under congruence.
\end{definition}

\begin{definition}[Extensional Equivalence]
\label{def:ext-equiv}
Two programs $f, g : A \to B$ are \emph{extensionally equivalent}
if $\forall x : A.\; f(x) = g(x)$.  In CTTT, this is
$\Path_{A \to B}(f, g)$, which by funext is equivalent to
$\prod_{x:A} \Path_B(f(x), g(x))$.
\end{definition}

\begin{theorem}[Structural/Extensional Boundary]
\label{thm:boundary}
Let $\mathcal{E}_{\mathsf{str}} \coloneqq \{(f, g) : f \sim_{\mathsf{str}} g\}$
be the set of structurally equivalent program pairs, and
$\mathcal{E}_{\mathsf{ext}} \coloneqq \{(f, g) : \forall x.\, f(x) = g(x)\}$
the set of extensionally equivalent pairs.  Then:
\begin{enumerate}
  \item $\mathcal{E}_{\mathsf{str}} \subsetneq \mathcal{E}_{\mathsf{ext}}$
    (structural equivalence is strictly weaker).
  \item The hard-8 pairs witness the gap:
    each is in $\mathcal{E}_{\mathsf{ext}} \setminus
    \mathcal{E}_{\mathsf{str}}$.
  \item D20 (specification uniqueness) bridges the gap: for any pair
    $(f, g) \in \mathcal{E}_{\mathsf{ext}}$ satisfying a
    \emph{recognizable} deterministic specification $S$, the path
    $\mathsf{det}_S$ provides a proof of equivalence.
  \item The D20 bridge is \emph{incomplete}: there exist extensionally
    equivalent programs that satisfy no specification expressible in
    the Z3 theory.
\end{enumerate}
\end{theorem}

\begin{proof}[Proof sketch]
\textbf{(1)} Every structural rewrite preserves extensional equality
(each axiom is sound), so
$\mathcal{E}_{\mathsf{str}} \subseteq \mathcal{E}_{\mathsf{ext}}$.
Strictness follows from (2).

\textbf{(2)} For each hard-8 pair, we verify:
(a) the pair is extensionally equivalent (both satisfy the same
deterministic specification, confirmed by testing on the benchmark
argument sets and by mathematical proof); and
(b) the bidirectional BFS in \texttt{search\_path} fails to find a
path within depth 4 (empirically verified), and the HIT structural
decomposition (\texttt{hit\_path\_equiv}) returns \texttt{None}.

\textbf{(3)} Given a deterministic specification $S$ with both $f$ and
$g$ satisfying $S$, Theorem~\ref{thm:det-spec} provides
$\Path(f, g)$.  The specification identification
(\texttt{\_identify\_spec}) automates recognition for common specs.

\textbf{(4)} Incompleteness follows from the fact that Z3's theory
(even with shared function symbols) cannot express all mathematical
properties.  For example, two programs computing the Ackermann function
via different recursion schemes are extensionally equivalent but have
no Z3-expressible specification connecting them (Z3 cannot reason
about non-primitive-recursive functions).
\end{proof}

\begin{remark}[The Three Strata]
The CTTT framework partitions program equivalences into three strata:
\begin{enumerate}
  \item \textbf{Structural stratum} ($\sim_{\mathsf{str}}$):
    Proved by syntactic rewriting.  Decidable (bounded BFS on finite
    axiom set).  Covers ${\sim}$42 of the 50~EQ benchmark pairs.
  \item \textbf{Specification stratum} (D20):
    Proved by recognizing a common deterministic specification.
    Semi-decidable (spec identification may fail, but if it succeeds,
    the proof is constructive).  Covers the hard-8 pairs.
  \item \textbf{Extensional stratum} (beyond D20):
    Proved only by mathematical argument outside the formal system.
    Not mechanizable in general (Rice's theorem).  No benchmark pairs
    fall here, but pathological programs do.
\end{enumerate}
\end{remark}

%% ═══════════════════════════════════════════════════════════
%% The Benchmark Pairs Table (updated + expanded)
%% ═══════════════════════════════════════════════════════════

\section{The Full Benchmark Classification}
\label{sec:benchmark-classification}

We classify all 50 equivalent pairs from the hard benchmark suite by
the proof technique required.  The structural pairs are resolved by
path constructors D1--D19 and D21--D24; the hard-8 pairs require D20.

\begin{center}
\small
\begin{tabular}{r l l l}
  \textbf{ID} & \textbf{Problem} & \textbf{Path} &
  \textbf{Key Axiom} \\
  \hline
  01 & Flatten + sort + unique & D4, D19 & Comprehension fusion, sort \\
  02 & Word frequency histogram & D13, D19 & Counter $\leftrightarrow$ fold \\
  03 & Arithmetic expression eval & D1 & Fix $\leftrightarrow$ fold \\
  04 & Register machine simulation & D1 & Fix $\leftrightarrow$ fold \\
  05 & Topological sort (lex-smallest) & D15 & Traversal order \\
  06 & RLE encode/decode roundtrip & D1, D4 & Fold + fusion \\
  \textbf{07} & \textbf{Convex hull} & \textbf{D20} &
    \textbf{Spec uniqueness} \\
  \textbf{08} & \textbf{Edit distance} & \textbf{D20} &
    \textbf{Spec uniqueness} \\
  09 & LRU cache simulation & D11 & ADT refinement \\
  10 & Matrix multiplication & D4, D5 & Fold fusion \\
  11 & Merge $k$ sorted iterators & D5 & Fold universality \\
  12 & Binary tree serialize/deserialize & D1 & Fix $\leftrightarrow$ fold \\
  13 & DAG path counting & D1, D16 & Fix + memo/DP \\
  14 & Calculator (parse + eval) & D1, D2 & Fix + inline \\
  15 & Longest increasing subseq & D16 & Memo $\leftrightarrow$ DP \\
  \textbf{16} & \textbf{Strongly connected comp} & \textbf{D20} &
    \textbf{Spec uniqueness} \\
  17 & Power set generation & D1 & Fix $\leftrightarrow$ fold \\
  18 & JSON-like pretty printer & D1, D2 & Fix + inline \\
  19 & 0/1 Knapsack & D16 & Memo $\leftrightarrow$ DP \\
  20 & Structural hashing (SHA-256) & D1, D19 & Fix + sort \\
  21 & Prime factorization & D1 & Fix $\leftrightarrow$ fold \\
  22 & Cycle detection (Floyd/Brent) & D1 & Fix $\leftrightarrow$ fold \\
  23 & Balanced parentheses & D9 & Stack $\leftrightarrow$ recursion \\
  24 & Levenshtein distance & D16 & Memo $\leftrightarrow$ DP \\
  25 & Merge overlapping intervals & D19 & Sort-then-scan \\
  26 & $k$-th smallest (quickselect) & D1 & Fix $\leftrightarrow$ fold \\
  27 & Longest common subsequence & D16 & Memo $\leftrightarrow$ DP \\
  28 & Regex state machine & D1, D9 & Fix + stack \\
  \textbf{29} & \textbf{Deep copy + sort} & \textbf{D20} &
    \textbf{Spec uniqueness} \\
  30 & Trie search & D11 & ADT refinement \\
  31 & Dijkstra shortest path & D10 & Priority queue \\
  32 & Island counting (grid) & D15 & BFS $\leftrightarrow$ DFS \\
  \textbf{33} & \textbf{$n$-th permutation} & \textbf{D20} &
    \textbf{Spec uniqueness} \\
  34 & Postfix expression eval & D9 & Stack $\leftrightarrow$ recursion \\
  35 & all $\leftrightarrow$ not-any-not & ALG & De Morgan involution \\
  36 & Transpose dict-of-lists & D4, D5 & Fold fusion \\
  \textbf{37} & \textbf{Fibonacci mod $p$} & \textbf{D20} &
    \textbf{Spec uniqueness} \\
  38 & Group anagrams & D19, D20 & Sort + spec \\
  39 & Perfect power check & D1 & Fix $\leftrightarrow$ fold \\
  \textbf{40} & \textbf{Binomial coefficient} & \textbf{D20} &
    \textbf{Spec uniqueness} \\
  41 & Zip-longest with fill & D5, D6 & Fold + lazy/eager \\
  42 & Currying & D24 & $\eta$-expansion \\
  43 & CSV parsing & D1, D4 & Fix + fusion \\
  44 & Coin change DP & D16 & Memo $\leftrightarrow$ DP \\
  45 & Grid path counting & D16 & Memo $\leftrightarrow$ DP \\
  46 & Mini type inference & D1, D8 & Fix + section merge \\
  \textbf{47} & \textbf{Stock profit} & \textbf{D20} &
    \textbf{Spec uniqueness} \\
  48 & Trailing zeroes in $n!$ & D1 & Fix $\leftrightarrow$ fold \\
  49 & Deep equality with cycles & D1, D9 & Fix + stack \\
  50 & $N$-queens & D1, D8 & Fix + section merge \\
\end{tabular}
\end{center}

Note: eq35 (all vs.\ any) is resolved by the De Morgan axioms
(boolean paths), not by specification reasoning.  The bold entries
are the hard-8 pairs requiring genuine D20 specification-level
arguments.

The two NEQ pairs at the algorithm boundary (neq06: late-binding vs.\
early-binding closure; neq20: in-place $+$$=$ vs.\ reassignment) are
correctly distinguished by the \emph{effect discipline}: mutation
effects and closure-capture timing use different Z3 symbols, so
the programs produce structurally different OTerms that cannot be
connected by any path.

